\section{Related Work}
\label{sec:related-work}

\subsection{Deepfake Detection}
\label{subsec:deepfake-detection}

Detecting synthetic images is a problem that has evolved alongside the generators themselves. Early work focused on face-swap forgeries~\cite{rossler2019faceforensics, perov2020deepfacelab}, where detectors exploited facial inconsistencies and GAN upsampling artifacts that were relatively easy to identify. As text-to-image diffusion models replaced GANs as the dominant synthesis paradigm, these detectors quickly became obsolete; the artifacts they relied on simply do not appear in diffusion-generated images~\cite{ojha2023universal}. More recent benchmarks like \textsc{OpenFake}~\cite{livernoche2025openfake} reflect this shift, covering a broad range of modern generators and content types far beyond the face-centric datasets that shaped early detection research.

\subsection{Diffusion-Based Image Generation}
\label{subsec:diffusion-generation}

Unlike GANs, which generate images in a single forward pass, diffusion models work by starting from pure random noise and gradually denoising it into a coherent image~\cite{ho2020ddpm}. Modern text-to-image systems build on this idea, accepting a text prompt as input and producing photorealistic images that no longer carry the upsampling artifacts GAN-era detectors were built to find.


\subsection{Detection Models}
\label{subsec:detection-models}

\subsubsection{CNNDetect (Wang et al., 2020)}
\label{subsubsec:cnndetect}
Wang et al.~\cite{wang2020cnndetect} proposed a CNN-based forensic detector for identifying synthetic images generated by convolutional neural network architectures. The work investigates whether images produced by CNN-based generators contain detectable artifacts that distinguish them from real images. Their detector is based on a ResNet-50 binary classifier trained to classify images as real or fake.

The original study mainly focused on CNN/GAN-based image generators, including ProGAN~\cite{karras2017progressive}, StyleGAN~\cite{karras2019style}, BigGAN~\cite{brock2018large}, CycleGAN~\cite{zhu2017unpaired}, StarGAN~\cite{choi2018stargan}, and GauGAN~\cite{park2019semantic}. These models rely on convolutional decoding and upsampling operations, which may introduce statistical artifacts in the generated image. To improve robustness, the authors also studied the effect of data augmentation, especially blur and JPEG compression.
This model is relevant to our work because it represents a classical GAN-era forensic baseline. However, modern diffusion models such as Midjourney and Stable Diffusion generate images through a different synthesis process. Therefore, evaluating Wang et al.'s detector on diffusion-generated images allows us to test whether CNN/GAN artifact-based detectors generalize to modern diffusion generators.

\subsubsection{EfficientNet (Tan and Le, 2019)}
\label{subsubsec:efficientnet}
EfficientNet~\cite{tan2019efficientnet} is a convolutional neural network family designed to achieve high accuracy with significantly improved efficiency. It is primarily characterized by a compound scaling method, which uniformly scales a model's depth, width, and resolution using a simple yet effective set of fixed coefficients.

This strategy is based on the intuition that higher-resolution images need more layers and more channels to capture fine-grained patterns. The process involves two steps: first, using neural architecture search to create an efficient baseline model (EfficientNet-B0), and second, performing a small grid search on that baseline to determine the best scaling constants for the entire family of models. Unlike the other detection models in our study, EfficientNet was not originally designed for deepfake detection; we include it as a general-purpose visual classification baseline to assess whether standard image classifiers can be repurposed for this task.

\subsubsection{CLIP+LoRA}
\label{subsubsec:clip-lora}

A key challenge in AI-generated image detection is that detectors trained directly for real-versus-fake classification may not learn a general notion of image authenticity. Instead, they can overfit to the artifacts of the generators present in the training set. Ojha et al. \cite{ojha2023universal} showed that such models may form an asymmetric decision boundary in which the fake class corresponds to the specific fingerprints learned during training, while the real class acts as a ``sink'' for samples that do not exhibit those fingerprints. This behavior limits generalization across generator families: artifacts that are shared among related models may support within-family transfer, but they do not necessarily transfer to other families of generators.

CLIP offers a different starting point for this problem. Radford et al. \cite{radford2021clip} introduced CLIP as a vision-language model trained on large-scale image-text pairs using contrastive learning. Its visual encoder and text encoder map images and natural-language descriptions into a shared embedding space, where matched image-text pairs are encouraged to be close and mismatched pairs are pushed apart. Since CLIP is trained for broad image-text alignment rather than for a real-versus-fake classification objective, its visual features are not directly specialized to the artifacts of a particular generator. This makes CLIP useful as a general representation space for fake-image detection.

Based on this idea, Ojha et al. \cite{ojha2023universal} used the frozen visual features of CLIP for fake-image detection instead of training a detector directly on generator artifacts. They showed that even simple classifiers, such as nearest-neighbor classification and linear probing, can generalize better to unseen generators. Cozzolino et al. \cite{cozzolino2024clipdetect} followed the same direction but built a stronger detector by training an SVM on CLIP features extracted from paired real and generated images. Both approaches keep the CLIP backbone frozen, which preserves its general-purpose representation. However, since CLIP itself is not adapted to the detection task, these methods can only use the real-versus-fake separation already present in the pretrained CLIP feature space.

A more adaptive alternative is to modify CLIP for the detection task, but full fine-tuning is costly and can overfit when only limited generated data are available. Low-Rank Adaptation (LoRA), introduced by Hu et al. \cite{hu2022lora}, addresses this trade-off by freezing the pretrained weights and learning low-rank updates in selected layers. Cappelletti et al. \cite{cappelletti2024lora} applied LoRA to CLIP for few-shot deepfake detection by inserting adapters into the visual backbone and training them together with a lightweight binary classifier. This allows the model to reshape CLIP's embedding space toward real-versus-fake separation while preserving much of its pretrained generalization capability.

These works motivate CLIP+LoRA as a suitable direction for diffusion-generated image detection. While frozen CLIP-based detectors preserve the generality of CLIP's pretrained visual representation, LoRA provides a lightweight mechanism for adapting this representation to the real-versus-fake task without fully fine-tuning the backbone. This is especially relevant for modern benchmarks such as OpenFake \cite{livernoche2025openfake}, which include realistic images from recent diffusion generators.


\subsubsection{DIRE (Wang et al., 2023)}
\label{subsubsec:dire}
Wang et al.~\cite{wang2023dire} introduced a novel image representation based on the concept of diffusion reconstruction error. The core theory is that if images are inverted and reconstructed using a pretrained diffusion model, the resulting reconstruction closely resembles the original in the case of diffusion-generated images. In contrast, real images are foreign to the diffusion model and result in a larger reconstruction error upon the same process.
The pipeline operates in two stages. First, each input image is inverted and reconstructed using a pretrained ADM (Ablated Diffusion Model). Second, the absolute pixel-wise difference between the original image and its reconstruction (the DIRE map) is computed and fed into a ResNet-50 classifier trained to distinguish real from fake based on this reconstruction error signal. 
The authors dedicated considerable effort to selecting the optimal pretrained diffusion model and classifier architecture for this pipeline, ultimately training on the LSUN-Bedroom dataset using an ADM-based reconstructor. Under these conditions, DIRE achieved state-of-the-art performance, reporting an average accuracy of 99.9\% and an average precision of 100\% across eight distinct diffusion-based generators, demonstrating strong generalization within the diffusion domain. However, since it was trained on this specific dataset, this makes it domain-specific by design.

\subsubsection{FIRE (Chen et al., CVPR 2025)}
\label{subsubsec:fire}
The Frequency-guIded Reconstruction Error (FIRE) model~\cite{chen2025fire} builds on the observation that diffusion models struggle to accurately reproduce mid-frequency information compared to real images.
The paper introduces a Frequency-based Mid-frequency Reconstruction Error (FMRE) module designed to identify mid-frequency filters. These filters are applied to create new images in which the mid-frequency information is removed. The resulting images are referred to as pseudo-generated images because they resemble diffusion-generated images, where mid-frequency information is typically absent.
The filtered images are then passed through a variational autoencoder (VAE) to project the image into the latent diffusion space and reconstruct it. Both the original and pseudo-generated images undergo this reconstruction process, after which the absolute difference between the reconstructed image and the original input image is computed.
The paper reports that the reconstruction error for real images, before and after filtering, is larger than that of diffusion-generated images, since diffusion-generated images are inherently more similar to the latent diffusion space.
FIRE differs from DIRE(Wang et al., 2023) in that it operates directly on raw RGB images at inference time, bypassing the computationally expensive per-image diffusion reconstruction step. This makes evaluation significantly faster while still leveraging frequency-domain features learned during training.